\documentclass[12pt,a4paper]{article}
\usepackage[utf8]{inputenc}
\usepackage[italian]{babel}
\usepackage{amsmath, amssymb, amsfonts, bm}
\usepackage{graphicx}
\usepackage{pgfplots}
\pgfplotsset{compat=1.18}
\usepackage{geometry}
\geometry{margin=2.5cm}
\usepackage{caption}
\usepackage{hyperref}

\title{\centering Studio matematico sulla coscienza umana \\ Legge 1: Legge del Tempo}
\author{Autore: Giuseppe Carlino}
\date{\today}

\begin{document}

\maketitle

\section{Premessa}

La \textbf{Legge del Tempo} descrive il modo in cui il tempo vissuto dall’individuo $t$ si integra con la conoscenza $\bm{k}$ accumulata (Legge 0) per influenzare la coscienza.  
Il tempo non è una semplice variabile lineare, ma uno stato dinamico soggetto a accelerazioni, decelerazioni e retroazioni dovute alla memoria e alla percezione della conoscenza.  

Questa legge funge da fondamento per le leggi successive: il tempo accumulato $t$ determina la crescita della conoscenza $\bm{k}(t)$, condiziona la capacità di scelta e influisce sulla volontà.

\section{Formalizzazione concettuale}

Si considera che il tempo percepito dalla coscienza $t_c$ possa differire dal tempo oggettivo $t$ in funzione della conoscenza:

\begin{equation}
\frac{dt_c}{dt} = 1 + \beta \frac{\|\bm{k}(t)\|}{\|\bm{k}_{\max}\|},
\end{equation}

dove:
\begin{itemize}
    \item $t$ è il tempo oggettivo misurabile;
    \item $t_c$ è il tempo soggettivo percepito dalla coscienza;
    \item $\bm{k}(t)$ è il vettore della conoscenza accumulata;
    \item $\bm{k}_{\max}$ è il limite massimo di conoscenza in ciascun ambito;
    \item $\beta$ regola quanto la percezione del tempo sia influenzata dalla conoscenza.
\end{itemize}

Questa equazione mostra che la percezione del tempo accelera con l’aumento della conoscenza: più un individuo sa, più “rapido” appare il tempo interno.

\subsection{Sistema dinamico integrato}

Combinando la Legge 0 (Conoscenza) e la Legge del Tempo, otteniamo un sistema dinamico:

\[
\begin{cases}
\displaystyle \frac{d \bm{k}}{dt} = \text{diag}(\bm{\alpha}(t))\, (\bm{k}_{\max}-\bm{k}) + \bm{\xi}(t),\\[1mm]
\displaystyle \frac{dt_c}{dt} = 1 + \beta \frac{\|\bm{k}(t)\|}{\|\bm{k}_{\max}\|}.
\end{cases}
\]

Questo sistema evidenzia la retroazione: la conoscenza influenza la percezione del tempo, e il tempo vissuto regola la crescita della conoscenza.

\section{Diagramma concettuale}

\textbf{Percezione del tempo in funzione della conoscenza:}

\begin{center}
\begin{tikzpicture}
\begin{axis}[
    width=14cm, height=8cm,
    xlabel={$t$ (tempo oggettivo)},
    ylabel={$t_c$ (tempo percepito)},
    grid=major,
    tick label style={font=\small},
]

% Curva di tempo percepito con perturbazioni simulate
\addplot[thick, blue, domain=0:10, samples=200] 
    {x + 3*(1-exp(-0.3*x)) + 0.2*rand}; 
\addlegendentry{$t_c(t)$ percezione tempo}

% Linea diagonale tempo lineare
\addplot[red, thick, dashed] coordinates {(0,0) (10,10)};
\addlegendentry{tempo lineare $t$}

\node at (axis cs:10,10.5) [anchor=south west] {tempo lineare};

\end{axis}
\end{tikzpicture}
\end{center}

\caption*{La curva blu mostra come la percezione del tempo $t_c$ accelera con la conoscenza accumulata $\bm{k}(t)$. La linea rossa tratteggiata rappresenta il tempo lineare oggettivo.}

\section{Interpretazione e implicazioni}

\begin{enumerate}
    \item Il tempo percepito dalla coscienza non è lineare: cresce in funzione della conoscenza accumulata.
    \item La legge introduce retroazioni: la conoscenza influenza il tempo percepito, creando un sistema dinamico interconnesso.
    \item Stabilisce la base temporale per la Legge della Scelta, poiché la capacità decisionale dipende dalla quantità di tempo percepito e dalla conoscenza.
    \item Perturbazioni $\bm{\xi}(t)$ e variazioni di $\bm{\alpha}(t)$ modellano la complessità della crescita cognitiva e temporale.
    \item Sistemi artificiali, come robot, non possiedono una percezione soggettiva del tempo: la loro "coscienza temporale" coincide con il tempo lineare $t$, mostrando un limite fondamentale rispetto agli esseri umani.
\end{enumerate}

\section{Relazione con le leggi successive}

\begin{itemize}
    \item La \textbf{Legge della Scelta} utilizzerà $t_c$ e $\bm{k}(t)$ per modellare la capacità decisionale.
    \item La \textbf{Legge della Volontà} dipenderà da $t_c$, $\bm{k}(t)$ e dalle scelte effettuate, rappresentando la capacità di agire.
    \item La \textbf{Legge Generale} (quadristato) integrerà tutti gli stati, mostrando come la coscienza umana sia un sistema complesso, multidimensionale e soggetto a retroazioni e perturbazioni.
\end{itemize}

\end{document}